% ============================================================
\chapter*{Pr\'eface}
\addcontentsline{toc}{chapter}{Pr\'eface}
% ============================================================

\section*{Objectifs du cours}

L'optimisation convexe occupe une place centrale dans les math\'ematiques appliqu\'ees
modernes. Elle fournit un cadre th\'eorique rigoureux et des algorithmes efficaces pour
r\'esoudre une vaste classe de probl\`emes apparaissant en apprentissage automatique,
traitement du signal, finance, recherche op\'erationnelle et bien d'autres domaines.

Ce cours de niveau Master s'adresse aux \'etudiants poss\'edant de solides bases en
analyse r\'eelle, alg\`ebre lin\'eaire et topologie. Il couvre les fondements th\'eoriques
de la convexit\'e, les conditions d'optimalit\'e, la dualit\'e, ainsi que les m\'ethodes
algorithmiques les plus importantes utilis\'ees en pratique.

L'approche adopt\'ee est \`a la fois rigoureuse sur le plan math\'ematique et tourn\'ee vers
les applications. Chaque concept th\'eorique est accompagn\'e d'interpr\'etations
g\'eom\'etriques, d'exemples concrets et d'impl\'ementations en Python.

\section*{Organisation du cours}

Le cours est organis\'e en dix chapitres, structur\'es selon une progression logique~:

\begin{enumerate}
    \item \textbf{Ensembles convexes} --- D\'efinitions fondamentales, propri\'et\'es
          topologiques, op\'erations pr\'eservant la convexit\'e, c\^ones, poly\`edres,
          th\'eor\`emes de s\'eparation.
    \item \textbf{Fonctions convexes} --- Caract\'erisations diff\'erentielles,
          in\'egalit\'es fondamentales (Jensen, gradient), continuit\'e et diff\'erentiabilit\'e.
    \item \textbf{Sous-diff\'erentiels et sous-gradients} --- Analyse convexe non lisse,
          calcul sous-diff\'erentiel, r\`egles de Moreau--Rockafellar.
    \item \textbf{Dualit\'e de Fenchel--Moreau--Rockafellar} --- Conjugaison convexe,
          th\'eor\`eme de biconjugaison, dualit\'e de Fenchel.
    \item \textbf{Conditions d'optimalit\'e --- KKT} --- Conditions n\'ecessaires et
          suffisantes, qualifications des contraintes.
    \item \textbf{Dualit\'e lagrangienne} --- Probl\`eme dual, saut de dualit\'e,
          dualit\'e forte, th\'eor\`eme de Slater, interpr\'etation \'economique.
    \item \textbf{Descente de gradient et variantes} --- Gradient \`a pas fixe et
          adaptatif, gradient acc\'el\'er\'e de Nesterov, gradient stochastique.
    \item \textbf{M\'ethodes proximales et ADMM} --- Op\'erateur proximal,
          algorithme ISTA/FISTA, ADMM, d\'ecomposition.
    \item \textbf{M\'ethodes de points int\'erieurs} --- Fonctions barri\`eres,
          m\'ethode du chemin central, complexit\'e polynomiale.
    \item \textbf{Applications} --- LASSO, SVM, r\'egression logistique,
          optimisation de portefeuille.
\end{enumerate}

\section*{Pr\'erequis}

Les pr\'erequis essentiels pour suivre ce cours sont~:

\begin{itemize}
    \item \textbf{Analyse r\'eelle} --- Suites, s\'eries, continuit\'e, diff\'erentiabilit\'e
          dans $\R^n$, th\'eor\`emes de convergence.
    \item \textbf{Alg\`ebre lin\'eaire} --- Espaces vectoriels, matrices, valeurs propres,
          d\'ecompositions matricielles (SVD, Cholesky), formes quadratiques.
    \item \textbf{Topologie} --- Ouverts, ferm\'es, compacit\'e, connexit\'e dans $\R^n$,
          notions d'int\'erieur relatif.
    \item \textbf{Programmation} --- Connaissance de base de Python (NumPy, matplotlib).
          La biblioth\`eque \texttt{cvxpy} sera utilis\'ee pour les impl\'ementations.
\end{itemize}

\section*{Notations}

Nous utilisons les notations suivantes tout au long du cours~:

\begin{center}
\renewcommand{\arraystretch}{1.4}
\begin{tabular}{cl}
    \toprule
    \textbf{Notation} & \textbf{Description} \\
    \midrule
    $\R^n$ & Espace euclidien de dimension $n$ \\
    $\inner{x}{y}$ & Produit scalaire $\sum_{i=1}^n x_i y_i$ \\
    $\norm{x}$ & Norme euclidienne $\sqrt{\inner{x}{x}}$ \\
    $\norm{x}_1$ & Norme $\ell_1$~: $\sum_i \abs{x_i}$ \\
    $\norm{x}_\infty$ & Norme $\ell_\infty$~: $\max_i \abs{x_i}$ \\
    $B(x,r)$ & Boule ouverte de centre $x$ et rayon $r$ \\
    $\overline{C}$ & Adh\'erence de l'ensemble $C$ \\
    $\mathrm{int}(C)$ & Int\'erieur de l'ensemble $C$ \\
    $\mathrm{ri}(C)$ & Int\'erieur relatif de $C$ \\
    $\mathrm{aff}(C)$ & Enveloppe affine de $C$ \\
    $\mathrm{conv}(S)$ & Enveloppe convexe de $S$ \\
    $\mathrm{cone}(S)$ & C\^one engendr\'e par $S$ \\
    $\mathrm{dom}(f)$ & Domaine effectif de $f$ \\
    $\mathrm{epi}(f)$ & \'Epigraphe de $f$ \\
    $f^*$ & Conjugu\'ee de Fenchel de $f$ \\
    $\partial f(x)$ & Sous-diff\'erentiel de $f$ en $x$ \\
    $\nabla f(x)$ & Gradient de $f$ en $x$ \\
    $\nabla^2 f(x)$ & Hessienne de $f$ en $x$ \\
    $\argmin_{x} f(x)$ & Ensemble des minimiseurs de $f$ \\
    $A \succeq 0$ & $A$ semi-d\'efinie positive \\
    $A \succ 0$ & $A$ d\'efinie positive \\
    $\mathbb{S}^n_+$ & C\^one des matrices SDP de taille $n$ \\
    $\iota_C$ & Fonction indicatrice convexe de $C$ \\
    $\sigma_C$ & Fonction d'appui de $C$ \\
    $\mathrm{prox}_f$ & Op\'erateur proximal de $f$ \\
    \bottomrule
\end{tabular}
\end{center}

\section*{Conventions}

\begin{itemize}
    \item Sauf mention contraire, tous les espaces consid\'er\'es sont de dimension finie
          sur $\R$.
    \item Les fonctions convexes consid\'er\'ees sont \`a valeurs dans
          $\R \cup \{+\infty\}$ (fonctions convexes propres).
    \item Le symbole $\square$ ou $\qed$ marque la fin d'une d\'emonstration.
    \item Les exercices sont class\'es par difficult\'e croissante~:
          $(\star)$ basique, $(\star\star)$ interm\'ediaire, $(\star\star\star)$ avanc\'e.
    \item Les blocs \texttt{algorithme} pr\'esentent le pseudo-code des m\'ethodes \'etudi\'ees.
    \item Les blocs \texttt{codeblock} contiennent du code Python ex\'ecutable.
\end{itemize}

\section*{Fil conducteur}

Le fil conducteur de ce cours est le probl\`eme d'optimisation convexe g\'en\'eral~:
\[
    \min_{x \in \R^n} \; f(x) \quad \text{sous} \quad
    g_i(x) \leq 0, \; i=1,\ldots,m, \quad
    h_j(x) = 0, \; j=1,\ldots,p,
\]
o\`u $f$ et les $g_i$ sont convexes et les $h_j$ sont affines. Nous \'etudions
successivement~:
\begin{itemize}
    \item la \emph{g\'eom\'etrie} de l'ensemble r\'ealisable (chapitres 1--2),
    \item l'\emph{analyse} des fonctions objectif (chapitres 2--4),
    \item les \emph{conditions d'optimalit\'e} et la \emph{dualit\'e} (chapitres 5--6),
    \item les \emph{algorithmes} de r\'esolution (chapitres 7--9),
    \item les \emph{applications} concr\`etes (chapitre 10).
\end{itemize}

\section*{Approche p\'edagogique}

Chaque chapitre suit une structure similaire~:
\begin{enumerate}
    \item \textbf{Motivation} --- Pourquoi ce sujet est important.
    \item \textbf{Th\'eorie} --- D\'efinitions, th\'eor\`emes avec d\'emonstrations compl\`etes.
    \item \textbf{Interpr\'etation g\'eom\'etrique} --- Illustrations TikZ et intuitions.
    \item \textbf{Exemples} --- Exemples d\'etaill\'es et contre-exemples.
    \item \textbf{Impl\'ementation} --- Code Python/cvxpy.
    \item \textbf{Exercices} --- Probl\`emes de difficult\'e vari\'ee.
\end{enumerate}

\section*{R\'ef\'erences principales}

\begin{enumerate}
    \item S.~Boyd et L.~Vandenberghe, \emph{Convex Optimization},
          Cambridge University Press, 2004.
    \item R.T.~Rockafellar, \emph{Convex Analysis},
          Princeton University Press, 1970.
    \item J.-B.~Hiriart-Urruty et C.~Lemar\'echal,
          \emph{Fundamentals of Convex Analysis}, Springer, 2001.
    \item Y.~Nesterov, \emph{Introductory Lectures on Convex Optimization},
          Springer, 2004.
    \item A.~Beck, \emph{First-Order Methods in Optimization}, SIAM, 2017.
    \item N.~Parikh et S.~Boyd, \emph{Proximal Algorithms},
          Foundations and Trends in Optimization, 2014.
    \item H.H.~Bauschke et P.L.~Combettes,
          \emph{Convex Analysis and Monotone Operator Theory in Hilbert Spaces},
          Springer, 2017.
    \item D.P.~Bertsekas, \emph{Convex Optimization Algorithms}, Athena Scientific, 2015.
\end{enumerate}

\vfill

\begin{flushright}
\textit{L'auteur}\\
Mars 2026
\end{flushright}
